\begin{helper}[Advance a repeated missing color by one path vertex]
Let \((\varphi,T)\) be a minimum bad pair with
\(T=(y_0,e_1,y_1,\ldots,e_n,y_n)\) and \(\rho(T)=1\).  Suppose
\(i+1<n\) and the color \(\beta\) is missing at both \(y_i\) and
\(y_n\).  Then there is a coloring \(\varphi'\) such that
\((\varphi',T)\) is again a minimum bad pair and \(\beta\) is missing
at both \(y_{i+1}\) and \(y_n\) under \(\varphi'\).
\end{helper}
\begin{proof}
Unfolding \noderef{TashTreeMinimumBadPair} and
\noderef{TashTreeBadPair}, and then \noderef{KTriple}, supplies
properness of \(\varphi\), the bound
\(\Delta(G)+1\leq k\), and the Tashkinov condition for \(T\).
There is a color \(\gamma\) missing at \(y_{i+1}\).  Indeed, for each
color in \noderef{PresentColors} at that vertex choose an exact edge
satisfying \noderef{Incident} there.  Two different present colors
necessarily have different chosen
edges, while two distinct incident edges with the same chosen color are
adjacent by \noderef{AdjacentEdges} and are forbidden by
\noderef{IsProperEdgeColoring}.  Hence the number of present colors is
at most \noderef{Degree} at \(y_{i+1}\), which is at most
\noderef{MaximumDegree} and therefore strictly below \(k\).  The
complement definition \noderef{MissingColors} supplies \(\gamma\).

The colors \(\beta,\gamma\) are distinct.  Otherwise both \(y_i\) and
\(y_{i+1}\) would miss \(\beta\).  These are distinct vertices by
injectivity in \noderef{TreeSequence}, and both lie in the proper prefix
through \(y_{i+1}\) because \(i+1<n\).  The elementariness supplied for
that prefix by \noderef{TashTreeProperPrefixElementary} would be
contradicted.

We next prove the two \noderef{TreeColorUnusedThrough} hypotheses
through \(y_{i+1}\).  Construct the prefixes used here with
\noderef{TreeSequencePrefix}.  If \(\gamma\) occurred on one of the displayed
colored edges \(e_2,\ldots,e_{i+1}\), the
\noderef{TashkinovTree} condition would provide an earlier displayed
vertex missing \(\gamma\).  That vertex is distinct from \(y_{i+1}\),
so the same elementary proper prefix would contain two vertices missing
\(\gamma\), impossible.  If \(\beta\) occurred on one of
\(e_2,\ldots,e_i\), its earlier Tashkinov witness would similarly be a
vertex distinct from \(y_i\) in the proper prefix through \(y_i\),
contradicting \noderef{TashTreeProperPrefixElementary}.  The remaining
edge \(e_{i+1}\) joins \(y_i\) to \(y_{i+1}\), since
\(\rho(T)=1\) means that \noderef{TashTreeSuffixPath} holds from the
first position.  This edge cannot have color \(\beta\), because it is
incident with \(y_i\), where \(\beta\) is missing by the definitions
\noderef{PresentColors} and \noderef{MissingColors}.  Thus neither
color is used through \(y_{i+1}\).

Apply \noderef{TashTreeControlledInterchange} to the distinct colors
\(\beta,\gamma\), the vertices \(y_i,y_{i+1}\), and these unused-color
facts.  It returns the exact component-swap coloring \(\varphi'\), says
that \((\varphi',T)\) is a minimum bad pair, and says that the original
color \(\beta\) is now missing at \(y_{i+1}\).  It remains to check its
missingness at \(y_n\).

The same controlled-interchange conclusion places \(y_i,y_{i+1}\) in
one two-color component and gives them subsingleton incidence neighbor
sets.  They are distinct.  Since \(y_n\) also misses \(\beta\),
\noderef{MissingColorIncidenceEndpoint} gives its vertex-node degree at
most one in the same two-color incidence graph.  If \(y_n\) lay in that
component, \noderef{TwoColorComponentSpanningPath} would give a simple
path spanning it with three distinct vertex-nodes of degree at most one:
\(y_i,y_{i+1},y_n\).  The component is nontrivial because it contains
the first two distinct vertices, so a simple path has exactly its two
ends as degree-at-most-one nodes.  This is impossible.  Hence \(y_n\)
is outside the swapped component.  The outside-vertex clause of
\noderef{TwoColorComponentSwapMissingBehavior} fixes its entire missing
set, so \(\beta\) remains missing there under \(\varphi'\).  This proves
all three asserted properties of the replacement coloring.
\end{proof}
